\chapter{Einleitung}

\section{Problemstellung und Motivation}

Phishing zählt zu den dauerhaft relevanten Bedrohungen der digitalen Wirtschaft. Allein im ersten
Quartal 2025 erfasste die Anti-Phishing Working Group mehr als eine Million Phishing-Angriffe;
die im selben Bericht wiedergegebene Branchenanalyse von OpSec Security ordnet
Online-Payment-Diensten sowie Finanz- und Bankdienstleistungen zusammen 30,9~\%
der von OpSec erfassten Angriffe zu (vgl. [APWG (2025), S. 3--4]).

Im Zentrum vieler Angriffe stehen manipulierte Webseiten, die vertrauenswürdige Dienste imitieren
und Nutzer zur Preisgabe schützenswerter Informationen verleiten (vgl. [APWG (2025), S. 2]).
Für ihre automatisierte Erkennung bilden insbesondere URL und HTML-Inhalt relevante
Informationsquellen. Während die URL Domain-, Pfad- und Parameterinformationen enthält, stellt
das HTML-Dokument textuelle, verlinkungsbezogene und strukturelle Eigenschaften der Webseite
bereit (vgl. [Dalton et al. (2026), S. 2--3]).

Die Erkennung wird dadurch erschwert, dass sich Phishing-Kampagnen und Webseiten fortlaufend
verändern. Generative Künstliche Intelligenz kann Angreifer zusätzlich bei der Erstellung
maßgeschneiderter Kampagnen unterstützen (vgl. [ENISA (2024), S. 67]). Gleichzeitig entstehen
regelmäßig neue Webseiten und Angriffsvarianten. Während Sperrlisten insbesondere bekannte
Phishing-Webseiten erfassen, können lernbasierte Verfahren URL- und HTML-Eigenschaften zur
Klassifikation bislang nicht beobachteter Webseiten nutzen (vgl. [Dalton et al. (2026), S. 1--3]). Eine hohe Leistung auf datensatznahen Testdaten erlaubt dabei noch keine belastbare Aussage über
bislang unbekannte Webseiten. Ähnliche Seiten oder gemeinsame Phishing-Kits können in Trainings-
und Testdaten auftreten und dadurch die Wiedererkennung bereits bekannter Strukturen begünstigen
(vgl. [Dalton et al. (2026), S. 3]).

Für einen betrieblichen Einsatz ist die Klassifikationsleistung nicht losgelöst von den hierfür
eingesetzten Ressourcen zu betrachten. Erstens ist relevant, welcher Anteil tatsächlicher
Phishing-Webseiten bei einer niedrigen False-Positive-Rate weiterhin erkannt wird, da Fehlalarme
legitime Webseiten betreffen und nachgelagerte Sicherheitsprozesse belasten können. Die
False-Positive-Rate wird entsprechend auch als Qualitätskennzahl von Security Operations Centers
gemeinsam mit Größen wie Eskalationen und Bearbeitungsaufwand betrachtet
(vgl. [ENISA (2020), S. 7]).

Zweitens stellt die Verfügbarkeit gelabelter Trainingsdaten eine Ressourcengrenze dar. Die
Beschaffung und zuverlässige Annotation hochwertiger Trainingsdaten verursacht zusätzlichen
Aufwand und gilt als wesentlicher Engpass datengetriebener Systeme
(vgl. [Roh et al. (2021), S. 1]). Drittens stellt sich bei einem leistungsstarken multimodalen
Klassifikationssystem die Frage, wie der zusätzliche Informationsumfang technisch effizient
genutzt werden kann. Für den Systementwurf ist damit relevant, ob eindeutige Fälle bereits auf
Basis der URL-Repräsentation abgeschlossen werden können, während unsichere Fälle weiterhin
die vollständige multimodale Verarbeitung durchlaufen. Vor dem Hintergrund begrenzter Budgets
und fehlender IT-Sicherheitsspezialisten, wie sie insbesondere für kleine und mittlere Unternehmen
berichtet werden, sind damit neben der maximalen Klassifikationsleistung auch der Bedarf an
gelabelter Supervision und die bedarfsgerechte Nutzung der verfügbaren Modellkomponenten
relevant (vgl. [ENISA (2021), S. 14]).

Self-Supervised Representation Learning (SSLR) setzt an dieser Ausgangslage auf Ebene der
Repräsentationsbildung an. Transformerbasierte Encoder können anhand ungelabelter Daten
selbstüberwacht vortrainiert und anschließend für eine gelabelte Downstream-Aufgabe genutzt
werden (vgl. [Devlin et al. (2019), S. 4171 und 4173--4175]). Darauf aufbauend kann das
Vortraining gezielt auf den ungelabelten Trainingsdaten einer konkreten Zielaufgabe fortgesetzt
werden. Gururangan et al. bezeichnen diese Weiteranpassung als Task-Adaptive Pretraining
(TAPT) und zeigen über mehrere Aufgaben hinweg Verbesserungen gegenüber ausschließlich
allgemein vortrainierten Modellen (vgl. [Gururangan et al. (2020), S. 8346--8347]). Für die Phishing-Webseitenerkennung ist damit zu untersuchen, welchen zusätzlichen Beitrag
TAPT gegenüber bereits allgemein vortrainierten Repräsentationen leistet, wie sich dieser Beitrag
bei unterschiedlicher Verfügbarkeit gelabelter Trainingsdaten verändert und wie die daraus
entstehende Repräsentationsbasis in ein leistungsstarkes multimodales Klassifikationssystem
überführt werden kann.

\section{Zielsetzung und Forschungsfragen}

Ziel der Arbeit ist die Konzeption, prototypische Entwicklung und empirische Evaluation eines
lernbasierten Klassifikationssystems zur Erkennung von Phishing-Webseiten anhand von URL-
und HTML-Daten. Neben der Klassifikationsleistung werden dabei der Umfang gelabelter
Downstream-Supervision und die Verarbeitungstiefe des resultierenden Systems betrachtet.

Zunächst wird untersucht, welchen zusätzlichen Beitrag eine task-adaptive selbstüberwachte
Anpassung gegenüber bereits allgemein vortrainierten Encodern leistet und wie sich dieser
Beitrag mit der verfügbaren gelabelten Supervision verändert. Darauf aufbauend wird geprüft,
ab welchem reduzierten Labelbudget das Leistungsniveau vollständig mit 200.000 Labels
trainierter Referenzverfahren erreicht wird. Eine klassische TF-IDF/XGBoost-Pipeline dient
ergänzend als Vergleichsverfahren; da sich dabei sowohl Repräsentation als auch Klassifikator
ändern, können Leistungsunterschiede nicht allein der task-adaptiven Anpassung zugeschrieben
werden.

Die Bewertung erfolgt auf dem vollständigen Testbestand sowie auf separat gebildeten domain-,
template- und zeitbezogenen Teilmengen an einem auf einer getrennten Kalibrationsmenge
bestimmten Low-FPR-Betriebspunkt. Anschließend wird die ausgewählte
Repräsentationskonfiguration um eine strukturelle DOM-Repräsentation und multimodale Fusion
zum TRI-Vollsystem erweitert. Darauf aufbauend wird untersucht, wie sich dieses System über
eine URL-first-Kaskade gestuft einsetzen lässt, wobei unsichere Fälle weiterhin vollständig
multimodal verarbeitet werden.

Aus dieser Zielsetzung ergibt sich folgende zentrale Forschungsfrage:

\begin{quote}
\textbf{Welchen Beitrag leistet Self-Supervised Representation Learning für ein lernbasiertes
Klassifikationssystem zur Erkennung bislang ungesehener Phishing-Webseiten und wie lässt sich
dieser Ansatz durch task-adaptive Repräsentationsanpassung sowie multimodale
Systemkomponenten weiterentwickeln?}
\end{quote}

Zur differenzierten Beantwortung werden vier Unterfragen formuliert:

\begin{description}
\item[FF1:] Wie verändert die task-adaptive Anpassung der URL- und HTML-Text-Repräsentationen
deren Downstream-Nutzbarkeit gegenüber nicht zusätzlich adaptierten vortrainierten Encodern?

\item[FF2:] Wie verändert sich der Zusatznutzen der task-adaptiven Repräsentationsanpassung bei
unterschiedlicher Verfügbarkeit gelabelter Trainingsdaten und unterschiedlichen
Downstream-Klassifikationsmechanismen?

\item[FF3:] Ab welchem reduzierten Labelbudget erreicht die task-adaptive RUT-Repräsentation
über die untersuchten Evaluationsbedingungen hinweg ein mindestens gleichwertiges
Leistungsniveau wie vollständig mit 200.000 Labels trainierte Referenzverfahren?

\item[FF4:] Wie verändern eine strukturelle DOM-Repräsentation, multimodale Fusion und eine
kaskadierte URL-first-Verarbeitung die Erkennungsleistung und den Inferenzaufwand des
prototypischen Klassifikationssystems?
\end{description}


\section{Vorgehensweise und Methodik der Arbeit}

Zur Beantwortung der Forschungsfragen wird ein empirisch-vergleichendes mit einem
konstruktionsorientierten Vorgehen verbunden. Als Datengrundlage dient der öffentlich verfügbare
PhreshPhish-Datensatz mit zeitlich verorteten URLs, gespeicherten HTML-Inhalten und
Klassenlabels für legitime und phishingbezogene Webseiten (vgl. [Dalton et al. (2026), S. 1--4]).

Zunächst werden allgemein vortrainierte Basisencoder mit Varianten verglichen, bei denen die
URL-Repräsentation, die HTML-Text-Repräsentation oder beide Informationszweige mittels TAPT
auf ungelabelten Webseitendaten weiter angepasst werden. Ihre Downstream-Nutzbarkeit wird mit
logistischer Regression und einem kleinen Mehrschicht-Perzeptron über sieben Labelbudgets
untersucht. Kontrastives Repräsentationslernen wird ergänzend anhand eines vorgelagerten
explorativen Versuchs eingeordnet.

FF3 vergleicht RUT über das gesamte Budgetraster mit R0@200k und TF-IDF/XGBoost@200k.
Für diese N10-Vergleiche werden Konfidenzintervalle, exakte Tests und korrigierte
Nichtunterlegenheitsentscheidungen berichtet. Modell- und Systementwicklung nutzen DEV,
die Schwellenbestimmung CAL und die Leistungsbewertung FINAL. Nachträglich ergänzte
Analysen werden gesondert gekennzeichnet. Eine externe Putra-Auswertung ergänzt die
Labelbudgetanalyse.

Anschließend wird die RUT-Basis schrittweise um HTML-Text, DOM und Gating erweitert.
Eine sequenzielle Komponentenablation prüft deren inkrementellen Beitrag. Das TRI-System
wird durch eine URL-first-Kaskade ergänzt, die nur unsichere Fälle vollständig multimodal
verarbeitet. Ein Frontend-Vergleich variiert die URL-Repräsentation zwischen Basis und TAPT.
Neben Klassifikationsmetriken werden Latenz, Durchsatz, GPU-Speicherbedarf und
Eskalationsrate erfasst.

Kapitel 2 und 3 begründen den Untersuchungsansatz; Kapitel 4 und 5 erläutern Methodik und
Umsetzung, Kapitel 6 die Ergebnisse und Kapitel 7 deren Einordnung.
